\begin{abstract}
\noindent
Physics-Informed Neural Networks (PINNs) provide a mesh-free framework for solving partial differential equations and are a natural candidate for option pricing, where the Black--Scholes equation reduces to a heat equation after a standard change of variables and propagates a non-smooth payoff backward in time into a smooth continuation-value surface. In this project, a fully connected neural network with parameters $\boldsymbol{w}$ is trained to approximate the option pricing function by minimizing a composite loss combining a PDE residual, a terminal payoff (initial) condition, and asset-price boundary conditions, with all derivatives computed by automatic differentiation. The framework is first validated on the Black--Scholes equation, where at the at-the-money point $S=K=100$, $\tau=0.5$ the PINN predicts a call price of $6.927$ versus the analytical value $6.889$, an absolute error of $0.0384$ (relative $0.56\%$) with a local PDE residual of $2.75\times 10^{-3}$. The method is then extended to the Heston stochastic-volatility model, where the network learns a function of $(S, v, \tau)$ and must satisfy a PDE containing a mixed second derivative; with both lower and upper $S$-boundaries imposed, the PINN reaches $6.807$ at the same point against a semi-analytical Heston benchmark of $6.837$, an absolute error of $0.0302$ (relative $0.44\%$). Finally, the same architecture is repurposed for inverse Heston calibration: the structural parameters $\kappa$, $\theta$, $\xi$, $\rho$ are made trainable and inferred jointly with $\boldsymbol{w}$ from 200 noisy synthetic option prices, using a two-stage Adam plus L-BFGS optimizer with a Feller-condition penalty. The inverse PINN attains an excellent price fit (mean absolute error $0.0286$, median relative error $\approx 0.02\%$, $R^2 \approx 0.999999$), but the recovered parameter vector $(\kappa, \theta, \xi, \rho) = (1.59,\, 0.0266,\, 0.104,\, 0.994)$ does not match the data-generating values $(2,\, 0.04,\, 0.3,\, -0.7)$; most strikingly, the correlation is recovered with the wrong sign---positive rather than negative---despite a near-perfect price reconstruction. This price/parameter dissociation, with $\rho$ entering the PDE only through the mixed-derivative term and therefore identifiable only through skew information that 200 sparsely placed prices cannot carry, is the central finding of the inverse experiment. The results suggest that PINNs are already strong differentiable forward pricing engines, but that inverse Heston calibration is genuinely ill-conditioned at this data density and will require richer collocation designs over a strike--maturity surface, economically motivated regularization of $\rho$ and $\xi$, or loss functions formulated in implied-volatility space rather than price space.
\end{abstract}
